\subsection{The complete good source-zero attachment}
\begin{statement}{P3K.83}{Proved}
Retain the actual source-zero family with
\((a,b)=(1,2)\), good-root order \(m=4p\), either realized norm
sign and every realized finite \(\delta\) character, with all source,
Cartier, root and annular maps prescribed above. In the selected
Tate fan \(T_0\), the positive fixed-scalar source loop has physical
integer index and oriented return class
\[
                         A_M=0,\qquad [\Sigma]=0\in\Theta_7.
\]
Its literal reflected fan \(JT_0\) has oriented return class
\(1\in\Theta_7\). The actual coupled loops
\(\lambda(v)=\lambda_0v^{km}\) have these same two classes,
\(0\) and \(1\), respectively. The reflected conclusion is an
oriented class; no additional integer characteristic pair is
assigned. The proof retains two chronological primitive maps,
their intermediate cap object, both charge-one faces, the entire
fixed trace and the actual physical Hopf completion. The already
proved vertical-fan recurrence is P3K.82.
\end{statement}
\begin{proof}
\par

\medskip\noindent\textbf{GC.1. The actual positive quotient and a proper weighted cut}

\par\nobreak

Use the full bilinear node coordinates of EA.42--46. The primary
associated resolution has coordinates (\(S_0\),\(S_1\),\(T_1\),z,V), with auxiliary
weights (1,1,-1,0,-1) and physical weights (0,0,0,1,-1). Its positive
auxiliary quotient is literally

\[
 W_1=\operatorname{Tot}\bigl(O(-1)_{T_1}\oplus O(-1)_V
                  \longrightarrow P^1_{[S_0:S_1]}\bigr)\times C_z.
 \tag{GC.1}
\]

The transition functions are the actual inverse powers of the two
homogeneous source coordinates. The source functions are
\(q_1\)=\(S_0\) z V and \(q_2\)=\(S_1\) \(T_1\). The original parameter v is the analytic
inverse of the base coordinate \(q_1\)-\(q_2\), with the unit retained as in
EA.17. This is the exact morphism from the original source, not
its homogeneous approximation.

Let H=\textbar{}\(S_0\)\textbar{} squared plus \textbar{}\(S_1\)\textbar{} squared, and define invariant functions

\[
 b=H|T_1|^2,\qquad c=H|V|^2,\qquad d=|z|^2,
 \qquad \rho^4=b+c^2+d^2,\qquad R^2=b+cd.          \tag{GC.2}
\]

The function \(\rho\) is proper on the vector-bundle neighbourhood of
the compact zero section. Every positive \(\rho\) level is smooth:
the derivative in \(T_1\) is nonzero when \(T_1\) is nonzero; otherwise a
nonzero V or z gives a nonzero derivative of its fourth power.
This is the weighted radial function with weights 2,1,1 on
(\(T_1\),V,z), respectively. R is exactly the source product radius
\textbar\textbar S\textbar\textbar{} \textbar\textbar T\textbar\textbar{} for T=(zV,\(T_1\)).

The physical fixed locus is exactly z=V=0. To see this without
discarding auxiliary stabilizers, the nonzero positive auxiliary
coordinates S force the auxiliary compensating phase to be one.
The remaining physical phase therefore fixes a point precisely
when both z and V vanish. This fixed locus is O(-1) over \(\mathbb P^1\), the
same exceptional trace disk as in RT.1. There is no isolated
positive-chamber fixed point in (GC.1). On z nonzero the function
z is an actual equivariant nonzero coordinate of physical weight
one; it supplies the physical line trivialization there.

Choose radii satisfying \(\epsilon\) positive and
\(r_1\)\^{}4\textgreater4 \(\epsilon\) squared, with \(r_1\)\textless{}\(\rho_0\). Retain the actual compact
domain

\[
 D_1=\{\rho\le r_1\}\ \cup
 \{r_1\le\rho\le\rho_0,\ R\ge\epsilon\}\ \cup
 \{r_1\le\rho\le\rho_0,\ d\ge c\}.              \tag{GC.3}
\]

Every union and intersection is in the same quotient (GC.1), by
the identity source, line and tangent map. At d=c one has
R squared=b+c squared and \(\rho\) fourth=b+2c squared, so
R squared is at least \(\rho\) fourth divided by two. Hence the
artificial d=c seam above \(r_1\) lies strictly in R\textgreater{}\(\epsilon\) and
supplies no boundary face. The entire zero section is in the
first piece. On the fixed trace c=d=0, R=\(\rho\) squared, so the
whole trace remains in (GC.3), including its exceptional disk.
The upper z arm is filled inside this same smooth quotient.

The removed horn is exactly \(r_1\)\textless{}\(\rho\)\textless=\(\rho_0\), c\textgreater d and R\textless{}\(\epsilon\).
It has V nonzero throughout. The two functions \(\rho\) and R have
regular transverse level faces there. Here is the check in the
actual complex fibre coordinates. At a point with \(T_1\) nonzero,
proportionality of their gradients would force the multiplier
between d(R squared) and d(\(\rho\) fourth) to be one. The nonzero
V coordinate would then give d=2c, contradicting c\textgreater d.~At a
point with \(T_1\) zero and R=\(\epsilon\) positive, both z and V are
nonzero. The ordered two-gradient determinant is
2(c squared-d squared)\textgreater0. If z is zero and \(T_1\) is nonzero,
the corresponding determinant is 2c\textgreater0. In the fully positive
case the \(T_1\) equation fixes the proportionality multiplier
to one, and the two others would force c=d=0. This handles zero
fibre coordinates without treating their squared norms as
smooth coordinates there. Invariant corner rounding retains
the specified outward normals.

\par

\medskip\noindent\textbf{GC.2. The complete charge-one lower shape}

\par\nobreak

In the actual V=1 chart there is no residual finite quotient.
The physical weights on (\(S_0\),\(S_1\),\(T_1\),z) are (-1,-1,+1,+1). Put

\[
 \sigma^2=\|S\|^2+|T_1|^2+|z|^2,\quad
 x=\|S\|^2/\sigma^2,\quad \zeta=|z|^2/\sigma^2.
\]

Direct substitution, with no radial coordinate map, gives

\[
 \rho^4=\sigma^4\bigl(x(1-\zeta)+\zeta^2\bigr),
 \qquad R^2=\sigma^4 x(1-x).                       \tag{GC.4}
\]

On the lower horn c\textgreater d says x\textgreater{}\(\zeta\). Its angular denominator
f0 squared=x(1-\(\zeta\))+\(\zeta\) squared is at most x. Consequently

\[
 1-x\le {R^2\over\rho^4}
       <{\epsilon^2\over r_1^4}<\frac14,
 \qquad 0\le\zeta\le1-x.                          \tag{GC.5}
\]

Thus x\textgreater3/4 and \(\zeta\)\textless1/4 on the closed horn. The square root f0 is therefore a positive
smooth angular function on a neighbourhood of that entire horn.
Choose an invariant cutoff equal to one there and supported in
x\textgreater1/2, and extend f=1+\(\chi\)(f0-1) to the whole charged sphere.
On that support \(\zeta\)\textless=1-x=x squared\textgreater1/4.
Together with the upper bound below this gives 1/2\textless f0\textless=1.
The cutoff support avoids every zero of f0; f is positive.
On the horn the exact defining-function equality is

\[
                         \rho^2=\sigma^2 f.         \tag{GC.6}
\]

Thus the complete outer charged shape \(\sigma\) squared times f
equal \(\rho_0\) squared agrees literally with the old radial face
on its lower rim. Since x+\(\zeta\)\textless=1, f0 squared is at most
(1-\(\zeta\)) squared plus \(\zeta\) squared, hence at most one. Our
cutoff therefore has f\textless=1. Take the round inner link with
\(\sigma\) squared equal to \(r_1\) squared divided by two. This lies
strictly below the entire horn, where \(\sigma\) squared\textgreater=\(\rho\) squared.
The interval region between that inner sphere and the full
outer shape is a compact free-circle annulus, with all source
functions carried by ambient inclusion. No radial map is
asserted to preserve \(q_1\) or \(q_2\).

Removing the horn from that charged annulus and gluing the
opposite resulting piece to D1 along its floor and side gives
the same explicit boundary-chain identity as EA.50--52. The
face maps here are the a=1 versions of EA.26 and EA.39, with
their actual node coefficient and scalar normalization. The
remaining complete charged boundary has physical line O(-1)
over \(\mathbb{CP}^3\): conjugating the two negative complex coordinates
preserves real eight-dimensional orientation, and its actual
coefficient map is {[}u,w{]} to w u in the tautological line.
Its physical pair is (-1,-4), so its characteristic index is
zero. This is the specified lower boundary calculation; it
does not count how the two primitive source steps occur in
the whole B strip.

The upper z region does not produce a charge-zero projective
substitution. It remains the actual region (GC.1), with the
equivariant nonzero function z on its free upper arm. Its line
and both boundary maps are retained in the ordinary source
gluing. The source functions satisfy
\textbar{}\(q_1\)\textbar+\textbar{}\(q_2\)\textbar\textless=R, so the same quantitative separation from the
outer source and parameter boundaries used in EA.19 applies.

\par

\medskip\noindent\textbf{GC.3. Two primitive cap steps with the full character and divisor}

\par\nobreak

Write m=4p. The actual primitive and suspension are

\[
 U_1:(p,w)\longmapsto(p-G,w/X(p-G)),\qquad
 B=-2\lambda q_1 U_1^2.                            \tag{GC.7}
\]

On the root \(q_{infty}\)=\(z_{infty}\)\^{}m the actual normalization from
EA.7 is

\[
 U_{\rm cap}=z_\infty^{-m/2}U_1,\quad
 F_0U_1=z_\infty^{m/2}U_{\rm cap}F_0,\quad
 q_1=z_\infty^{-m}h,\quad
 F_0B=-2\lambda h\,U_{\rm cap}^2F_0.              \tag{GC.8}
\]

All source and line factors of \(F_0\) are retained. The root map
Ucap is the literal translation with target coefficient
1/X'(p-G'), not an uncorrected tensor frame. The canonical
deck sends X' to -X', and the exact finite relation is
Ucap \(A_0\) Ucap inverse=-\(A_0\). Hence a single primitive step changes
the finite physical character by m/2; it is not declared a
self-map of the unchanged cap data.

Here is the complete intermediate two-step groupoid. For
\(\epsilon\)=0,1,2 set

\[
 A_\epsilon=(-1)^\epsilon A_0,\qquad
 F_\epsilon=z_\infty^{-\epsilon m/2}F_0.
 \tag{GC.9}
\]

Then the actual equalities are

\[
 U_{\rm cap} A_\epsilon U_{\rm cap}^{-1}
       =A_{\epsilon+1},\qquad
 F_{\epsilon+1}U_1=U_{\rm cap}F_\epsilon.
 \tag{GC.10}
\]

The first equality retains the full root generator, including
its source translation and chosen finite \(\delta\) character.
The scalar \(z_{infty}\) to the m/2 power has root character -1,
which cancels exactly that of Ucap on the ordinary overlap.
After two steps \(A_2\)=\(A_0\), but \(F_2\)=\(z_{infty}\) inverse-m times \(F_0\).
This vertical divisor is not removed as a unit. It is
compensated only by the actual \(q_1\) factor in B, giving the last
square in (GC.8), with h still present. Both norm signs n and
every original finite character remain throughout.

\par

\medskip\noindent\textbf{GC.4. The retained derivative and spin maps of the two steps}

\par\nobreak

For one graph reset EA.18--20 gives the full complex-four
determinant \(z_{s}\) inverse-squared in this branch. Its chosen
determinant root is the literal function \(z_{s}\) inverse. Thus two
consecutive reset charts carry the product
(\(z_{s}\),0 \(z_{s}\),1) inverse as their determinant root, evaluated at
the actual consecutive source points. Neither source coordinate
is replaced by its value on a central component. On another
projective source chart the canonical volume and this root are
transported by that chart's actual transition; these are bundle
arrows, not an assertion that \(z_{s}\) inverse extends as a function
through zero.

The charged overlap EA.26 has complex determinant -1 and the
fixed root i before the normalized charged gauges; it is the
identity on their ordered complex-four coefficients. For a=1
there is no residual mu2 character at this lower link. The
canonical complex spin-c lifts are the exterior-algebra lifts
of these actual derivatives. Their compositions are the lift
of the displayed composed coordinate arrow, with the stated
determinant roots. The positive norm changes have the exact
diagonal and nilpotent derivative homotopies of RT.15, based
at these arrows.

The vertical factor \(z_{infty}\) to the m/2 power in (GC.8) is also
kept as a typed root-line map. In an ordinary nonzero line
coefficient w its derivative has determinant \(z_{infty}\) to the
m/2 power: the source coordinates are fixed and the additional
base derivative is triangular. Since m/2=2p, the corresponding
determinant root is the single-valued \(z_{infty}\) to the p power
on this annulus. This statement is in the dw coefficient frame.
In the logarithmic line frame dw/w the same map has determinant
one and the extra derivative is unipotent; the frame arrow
dw=w dlog(w) is retained. In an associated node pair (U,V)
the simultaneous map is (\(z_{infty}\)\^{}(m/2) U,\(z_{infty}\)\^{}(-m/2) V),
whose two normal determinants cancel. None of these different
frames is equated without its displayed coefficient arrow.
At the root centre the scalar powers are maps of the
corresponding divisor lines, not nonvanishing scalar units.

The canonical complex spin-c lift of the one-line character
-1 has the literal representative {[}e1 e2,i{]}, using the
positive half-rotation in that complex line and \(e_{i}\) squared
equal to -1. Its square is {[}-1,-1{]}, the identity in Spin-c.
On the exterior-algebra spinors the first character acts as
plus one in degree zero and minus one in degree one. Thus the
two character changes in (GC.10) compose to the original
character with its canonical lift, not merely as bare signs
on the underlying line. The determinant frames of the source
reset are still the chronological roots stated above. The full
second-step norm is exactly the B square in (GC.8), so no
additional independent sign or unrecorded ruling reset is
inserted at the cap.

Sections GC.5--9 place this entire two-object sequence in the
actual finite source strip, matching the upper z arm, the two
local charged collars and the whole fixed trace. The displayed
local/cap composition is retained in that full assembly.

\par

\medskip\noindent\textbf{GC.5. The entire upper free product, with its coefficient arrow}

\par\nobreak

The a=1 upper region z nonzero has a literal product map that
also retains its R=0 exceptional curve. Set W=zV. In the actual
positive auxiliary quotient the map and its inverse are

\[
 (S,T_1,V,z)\longmapsto(S,T_1,W=zV,z),\qquad
 (S,T_1,W,z)\longmapsto(S,T_1,W/z,z).
 \tag{GC.11}
\]

Since z has auxiliary weight zero, these maps commute with every
auxiliary chart transition. Thus this region is exactly

\[
 \operatorname{Tot}\bigl(O(-1)_{T_1}\oplus O(-1)_W
          \longrightarrow P^1_S\bigr)\times C^*_z.
 \tag{GC.12}
\]

The physical circle acts only on the z factor. The full source
functions are \(q_1\)=\(S_0\) W and \(q_2\)=\(S_1\) \(T_1\), independent of that factor;
the original source parameter is carried by the same analytic
inverse of \(q_1\)-\(q_2\) as in GC.1. In particular \(T_1\)=W=0 is the entire
exceptional \(\mathbb P^1\) times C*, carried by (GC.11), not a point removed
from its domain. Its R=0 locus needs no positive-norm flow to
make this product map defined.

The physical line on this piece has a specified trivialization.
Use the convention that a physical positive circle frame p and
its coefficient w represent the vector wp; consequently
(gp,g inverse w) represents that same vector. The actual
weight-one coordinate z gives

\[
                  [p,w]\longmapsto z(p)w.             \tag{GC.13}
\]

The formula is invariant and invertible because z is nowhere
zero on this piece. It is the coefficient arrow in this chart,
not an assertion about its extension through z=0 or through the
next source reset. The complex derivative of (GC.11) in the
ordered (V,z) block is

\[
              \begin{pmatrix}z&V\\0&1\end{pmatrix},
 \tag{GC.14}
\]

with all source coordinates retained. This carries the canonical
complex spin-c structure by its actual exterior-algebra map.
Its determinant z is retained; no single-valued square root on
the z circle is asserted. The chart's physical line trivialization
alone does not discard a transition contribution at its boundary.

The marked G divisor in the final node coordinates is \(S_1\)-\(S_0\)=0,
by EA.45. In particular the coordinate \(S_0\)=F+ in (GC.11) is not
the old Cartesian coefficient s0. Any next-source transition
uses the actual EA.43--46 unit and the chronological reset,
before one assigns a marked source section to an upper point.
The full two-step outside attachment remains the next map to
be checked; (GC.11)--(GC.14) establish the complete upper chart
without assigning that attachment or an individual return.

\par

\medskip\noindent\textbf{GC.6. The proper upper graph and its actual primitive endpoint}

\par\nobreak

Put \(T_0\)=W=zV in (GC.11), and use the four invariant conifold
functions

\[
 q_1=S_0T_0,\quad q_2=S_1T_1,\quad
 C=S_0T_1,\quad D=S_1T_0,\qquad q_1q_2=CD.
 \tag{GC.15}
\]

The rational opposite source ruling {[}\(T_0\):\(T_1\){]} is retained by its
actual graph, not made into a coordinate at \(T_0\)=\(T_1\)=0. Its proper
common resolution has homogeneous coordinates {[}\(S_0\):\(S_1\){]}, {[}t0:t1{]},
a coefficient \(\eta\) of O(-1,-1), and the nonzero z. Its two
projections are

\[
 ([S],[t],\eta,z)\longmapsto
       ([S],T_0=\eta t_0,T_1=\eta t_1,z),
\]
\[
 ([S],[t],\eta,z)\longmapsto
       ([t],A_0=\eta S_0,A_1=\eta S_1,z).
 \tag{GC.16}
\]

Explicitly the two auxiliary scalings multiply S and t by
\(\alpha\) and \(\beta\) and \(\eta\) by (\(\alpha\) \(\beta\)) inverse. This gives
both inverse ordinary blowup charts wherever t0 or t1 is
nonzero. The first projection is precisely the blowup of the
normal ideal (\(T_0\),\(T_1\)) of the upper exceptional \(\mathbb P^1\); the second
contracts the other projective factor. Both retain the source
functions \(q_1\)=\(\eta\) \(S_0\) t0 and \(q_2\)=\(\eta\) \(S_1\) t1. The physical circle
acts only on z, so the graph and both proper source products
are free on this entire upper piece. Their physical coefficient
arrow is the pullback of (GC.13), including \(\eta\)=0.

Here the full \(\theta\) unit has a useful exact extension. Retain
Q=\(q_1\) \(q_2\) uQ, P=uP/\(q_2\) and uG=uQ uP. In the \(S_1\) chart TL uses
S=\(S_0\)/\(S_1\) and

\[
 E=\frac{2\Theta_Q(S)}{(1-S)\Theta_Q(S/P)}.
 \tag{GC.17}
\]

This unit descends to the full conifold (GC.15). Indeed its
normally convergent product factors are expressed by

\[
 QS=u_Qq_1C,\quad Q/S=u_Qq_2D,\quad
 S/P=C/u_P,\quad Q/(S/P)=u_GD.
 \tag{GC.18}
\]

For every higher power of Q multiply the first two expressions
by that power minus one. The apparent marked-section factor
1-S is already divided out. All numerator and denominator
factors are therefore holomorphic and nonzero on a sufficiently
small full conifold neighbourhood, including both S charts and
both target t charts. Their product converges normally there.
The value at the vertex is exactly two. This supplies an
extending logarithm of E/2 on that neighbourhood, without
replacing E by its vertex value. In particular

\[
                              L=E z
 \tag{GC.19}
\]

is the same actual nonzero physical coefficient over the whole
upper graph. It descends through each projection in (GC.16),
since E is a function of the four common invariant functions.

The endpoint is the actual inverse primitive. On the original
Tate torus, with the precise source/Cartier map of TL and SZ.43,
its input coefficients are

\[
              Z_0=u_GD,\qquad V_0=P(C/u_P)L.
 \tag{GC.20}
\]

The identity is first read in the \(S_1\) chart, where it is TL.11
with the actual U=\(S_0\) z coefficient. The homogeneous line-frame
arrow on the other S chart is exactly TL's retained Cartier
transition; (GC.18) gives its extension. Applying \(U_1\) inverse,
which has the literal torus map (Z,V) to (PZ,\(c_{X}\) V), gives

\[
 Z_1=Pu_GD,\qquad V_1=u_X C L,
                  \qquad c_X=q_2u_X .
 \tag{GC.21}
\]

No \(q_2\) factor or source translation has been discarded. On the
t1 chart of (GC.16), with t=t0/t1 and coordinates (\(A_0\),\(A_1\)),
these formulas read

\[
 q_1=A_0t,\quad q_2=A_1,\quad C=A_0,\quad D=A_1t,
\]
\[
 Z_1=u_Gu_Pt,\qquad V_1=u_XA_0L,\qquad
 L_1=Z_1V_1/q_1=(u_Gu_Pu_X)L .
 \tag{GC.22}
\]

The last expression is interpreted in its actual toric edge
chart. The unit character on the unshifted negative diagonal
\{e1,e2\} is ZV/\(q_1\). Thus \(L_1\) is an extending nonzero coefficient
at \(A_0\)=0 and t=0, rather than a division performed only after
restricting to that divisor. The source node coordinate
\(q_1\)/\(Z_1\) is \(A_0\)/(uG uP), with its full unit.

On the t0 chart put p=t1/t0, so \(q_1\)=\(A_0\), \(q_2\)=\(A_1\) p,
C=\(A_0\) p and D=\(A_1\). Carry the actual source period
h:(Z,V) to (QZ,PV), including P. Then

\[
 Z_2=QZ_1=q_1u_G^2D,\qquad
 V_2=PV_1=u_Pu_X(A_0/A_1)L,
\]
\[
             L_2=Z_2V_2/q_1^2=u_G L_1.
 \tag{GC.23}
\]

These are the other indexed node/edge coordinates. Relative to
the first \(q_1\) fan chart the h arrow translates the edge by e1;
its unit character is ZV/\(q_1\) squared. The full \(q_2\)-dependent
period P is still in (GC.23). The target edge coordinates on this second chart must also retain
the two S charts. With \(S_1\)=t0=1, \(\eta\)=r, \(S_0\)=s and t1=p, they are

\[
 q_1=rs,\quad q_2=rp,\quad
 Z_2/q_1=u_G^2r,\quad q_1^2/Z_2=s/u_G^2.
 \tag{GC.23a}
\]

All four quantities extend through r=s=p=0. With \(S_0\)=t0=1,
\(\eta\)=r, \(S_1\)=s and t1=p, one instead has \(q_1\)=r,\(q_2\)=rsp. The
ratio \(q_1\) squared over \(Z_2\) is not a coordinate at s=0. At that
endpoint retain the actual adjacent source Tate-node pair

\[
 w=Z_1^{-1}=p/(u_Gu_P),\qquad
 z_T=QZ_1=u_G^2r^2s,\qquad z_Tw=Q.
 \tag{GC.23b}
\]

The corresponding nonzero line coefficient is
Y0=\(Z_1\) \(V_1\)/\(q_1\)=\(L_1\). Its h-side frame is Yh=\(q_1\) uG Y0.
The latter factor is retained as a divisor, not called a unit:
the adjacent associated node coordinates are

\[
           (U_h,V_h)=(q_1u_GY_0,(u_GY_0)^{-1}),
                         \qquad U_hV_h=q_1.
 \tag{GC.23c}
\]

These are regular even at r=s=p=0 because Y0 and uG are
nonzero. The proper graph projection in (GC.16), with its
{[}S{]} and {[}t{]} coordinates, is still retained. No inverse to its
exceptional contraction is asserted. In the remaining two
charts t1=1 one has \(Z_1\)=uG uP t0 and
\(q_1\)/\(Z_1\)=\(\eta\) \(S_0\)/(uG uP), with \(q_2\)=\(\eta\) \(S_1\); these are regular for
either \(S_0\)=1 or \(S_1\)=1. Thus every endpoint of \(\mathbb P^1_S\times\mathbb P^1_t\)
has a specified source-node and associated-line-node map, not
an unextended ratio. On the common edge overlap \(L_2\)=uG \(L_1\) is
an extending unit, whose logarithm is carried with the original
Cartier map. This is not a free change of suspension marking.

The inverse primitive belongs to the two-object cap sequence
(GC.9)--(GC.10). In the direction used above its root square is

\[
                    F_\epsilon U_1^{-1}
                         =U_{\rm cap}^{-1}F_{\epsilon+1}.
 \tag{GC.24}
\]

Both finite characters are retained. The complete B inverse is
\(U_1\) inverse twice followed by the scalar (minus2 \(\lambda\) \(q_1\))
inverse, and its cap square is the inverse of (GC.8).
Equations (GC.16)--(GC.24) give the full upper proper graph and
its one-primitive endpoint through the intermediate object.
They do not replace that composite by its normal-cone limit or
assert that one primitive is an automorphism of the original
cap object. Section GC.9 glues the two chronological primitive
graphs through this intermediate object and the single whole
fixed trace, using the explicit trace of GC.8.

\par

\medskip\noindent\textbf{GC.7. The full two-step upper line character}

\par\nobreak

For the complete good suspension, put \(c_{lambda}\)=minus2 \(\lambda\).
The actual inverse, including its selected \(q_1\) factor, is

\[
 B^{-1}:(Z,V)\longmapsto
       (P^2Z,c_X^2V/(c_\lambda q_1)).
 \tag{GC.25}
\]

Substituting (GC.20), with CD=q1q2, gives

\[
 Z_B=P^2u_GD,\qquad
 V_B={q_2u_X^2C\over c_\lambda q_1}L,\qquad
 Z_BV_B={u_Gu_P^2u_X^2\over c_\lambda}L.
 \tag{GC.26}
\]

The product in the last expression is the unit character of
the actual height-zero indexed edge. Its coefficient is the
full extending base unit \(F_{B}\)=uG uP squared uX squared divided
by \(c_{lambda}\). \(F_{B}\)(0)=1/\(c_{lambda}\), and \(F_{B}\)/\(F_{B}\)(0) has an
extending logarithm. Both primitive source translations and
the \(q_1\) factor were used before this cancellation. In particular
(GC.26) does not discard the p-dependent multiplier in SZ.34.

The source graph of the two chronological primitive maps may
be taken with both projective intermediate source coordinates.
Its line factor is the same C* coordinate L, with its target
coefficient changed by the explicit nonzero \(F_{B}\). Thus the
actual upper source/line graph is the product of that proper
source graph with C*\_L under the typed target map
\(Y_{B}\)=\(F_{B}\) L. The physical action is scalar on this factor.
All source functions are pulled back by the graph projections.
This proves the physical-line trivialization on that graph,
including its projective source endpoints; it does not assert
that the source graph itself is smooth or that either of its
projections is an isomorphism at a contracted curve.

The full finite-root square is

\[
 F_0B^{-1}=(c_\lambda h)^{-1}U_{\rm cap}^{-2}F_0,
                    \qquad q_1=z_\infty^{-m}h.
 \tag{GC.27}
\]

It is the composition of the two different cap-object maps
in (GC.24), with \(F_2\)=\(z_{infty}\) inverse-m times \(F_0\) still retained.
The scalar character returns only after that composition.
No additional primitive ruling change is inserted at the cap.
The finite smooth cut of this full proper source graph is
constructed in GC.8--9, with its physical line and the whole
fixed trace carried on every collar.

\par

\medskip\noindent\textbf{GC.8. A finite smooth trace for the upper proper graph}

\par\nobreak

The upper source graph in (GC.16) has the following explicit
smooth real trace. Use its common source Cox coordinates
(s0,s1,t0,t1), with compact auxiliary circle weights
(+1,+1,-1,-1), and put

\[
 \mu=(\|t\|^2-\|s\|^2)/2,\qquad R=\|s\|\|t\|,
                  \qquad \rho_s^2=\|s\|^2+\|t\|^2.
 \tag{GC.28}
\]

Choose positive constants with r squared divided by two smaller
than both \(\delta\) and R0. Take the compact region

\[
 -\delta\le\mu\le\delta,\qquad R\le R_0,
                         \qquad \rho_s^2\ge r^2
 \tag{GC.29}
\]

and divide by that compact auxiliary circle. Compactness follows
from the exact identity \(\rho_{s}\) fourth=4mu squared+4R squared.
The auxiliary action is free off the removed origin because all
four nonzero coordinate weights have absolute value one. The
moment levels have nonzero differential away from the origin;
on the R=R0 side both s and t are nonzero, and the two moment
and product-radius differentials are independent. The round
inner boundary lies strictly between the two moment faces and
strictly inside the product-radius side by the chosen inequalities.
Invariant rounding therefore gives a compact smooth seven-manifold
with precisely its two moment faces, the stable side, and one
full inner quotient of \(S^7\) by this auxiliary circle.

The two moment faces are the actual two source small resolutions.
On the s-stable face, positive auxiliary normalization takes a
Cox orbit to that specified moment value and retains {[}s{]} and the
products \(s_{i}\) \(t_{j}\); on the t-stable face it retains {[}t{]} and those
same products. On their R\textgreater0 graph collar both projective factors
are present and the projections are exactly (GC.16). Each moment
face also includes its own small-resolution zero section: t=0
on the s-stable face and s=0 on the t-stable face. The opposite
projective factor is not asserted to be defined at that zero
section. These endpoints are supplied by the actual moment
face, while the proper graph retains both factors only on its
specified comparison collar. The positive auxiliary flow
(s,t) to (a s,a inverse t) has strictly monotone moment on each
nonzero orbit. On the R=R0 side it supplies the unique map to
any prescribed intermediate moment. It preserves every \(s_{i}\) \(t_{j}\),
and hence the full \(q_1\),\(q_2\), the analytic inverse giving v, and the
unit E in (GC.18). This is the typed stable-side identification,
not a radial normalization of q or v.

Carry the physical coefficient L of (GC.19) and take its unit
circle with respect to the actual transported metric. The physical
circle acts on this factor alone. The resulting eight-dimensional
trace carries the same physical line with the coefficient arrow
{[}p,w{]} to L(p)w. At a target face use its actual coefficient
\(L_1\) or \(L_2\) from (GC.22)--(GC.23); the metric is transported by
that full nonzero multiplier. The extending logarithms of the
unit factors give the explicit interpolation of coefficient
frames through the trace. Constants and their phases are kept.
For the composed endpoint use the full \(F_{B}\) in (GC.26).
Thus all maps on the stable side agree with the ordinary upper
finite strip, including both source- and line-node frame changes
of (GC.23a)--(GC.23c).

The one new full inner face is (\(S^7\)/\(S_1\)\_aux) times \(S_1\)\_L. Conjugating
the two negative auxiliary coordinates identifies its first
factor with \(\mathbb{CP}^3\) by the literal positive Hopf map; the physical
circle remains the independent second factor. Denote its oriented
physical quotient by \(Q_{up}\), with the orientation induced by this
actual boundary face. Its physical associated line has the global
nonzero frame just displayed. Consequently its first Chern class
is zero and both x cubed and \(p_1\) times x have zero integral.
No physical charge-one Hopf line is substituted for this trivial
line. In particular no topology of \(Q_{up}\) is inferred from its
vanishing characteristic numbers.

The trace's stable tangent and determinant maps are the actual
auxiliary quotient maps of these Cox coordinates. The sum of its
four auxiliary weights is zero, so the determinant frame is the
one descended from their ordered complex volume, with the orbit
and moment directions explicitly removed. The physical L circle
uses its original complex coefficient frame dL and its radial
normal. A passage to dL/L carries the factor L, rather than
silently changing the spin along that circle. Unit coefficient
changes are lifted with their actual extending logarithms; the
finite character maps are the canonical lifts of GC.4. This fixes
the determinant/spin-c maps on every interface of this free trace.
It creates no additional physical fixed stratum.

\par

\medskip\noindent\textbf{GC.9. Chronological two-primitive finite strip and its boundary ledger}

\par\nobreak

Fix the scalar by the actual RT.12 based homotopy. Put
\(c_{lambda}\)=minus2 \(\lambda\) and choose a nonzero constant c0. Factor

\[
 K_0=c_0U_1,\qquad K_1=(c_\lambda/c_0)q_1U_1,
                              \qquad B=K_1K_0.
 \tag{GC.30}
\]

The intermediate object is exactly (\(A_1\),\(F_1\)) from (GC.9). The
finite-root arrows on these two pieces are

\[
 F_1K_0=c_0U_{\rm cap}F_0,\qquad
 F_0K_1=(c_\lambda h/c_0)U_{\rm cap}F_1.
 \tag{GC.31}
\]

The second equality uses \(F_0\)=\(z_{infty}\) to the m times \(F_2\) and the
retained \(q_1\)=\(z_{infty}\) inverse-m times h. Thus the first root
character is shifted by m/2 and the second returns it, with its
canonical spin-c composition. Their product is exactly the
original B square, including h, norm sign and every finite
character. The primitive object is never identified with the
original one merely by its underlying source surface.

On the compact source and root complement of the collision graph,
choose c0 small enough to make the first transported norm cut
contracting. Then choose the fixed \(\lambda\) small enough that the
second cut is contracting too. Both maps and their inverses are
bounded there. Changing \(\lambda\) within this same nonzero range
uses the zero-winding radial path already retained in RT.12.
On the graph collars use the actual Fubini--Study norms
Nt=\textbar xi\textbar{} \textbar\textbar t\textbar\textbar{} and Ns=\textbar xi\textbar/\textbar\textbar s\textbar\textbar{} from EA.12. Their ratio R is
retained, and the actual positive norm map is the a=1 instance
of EA.13. Choose the cut between \(K_0\) and \(K_1\) as a distinct marked
face. This gives two chronological compact primitive pieces,
with one intermediate ordinary face and the root object (\(A_1\),\(F_1\)).
Their intermediate faces are glued by the actual coefficient
identity, with opposite outward normals. Their full finite-root
faces are glued by (GC.31), not by discarding a finite character.

There is one pure nonzero-line Cox-origin cut for \(K_0\). In its
common xi=1 gauge the source and target finite faces are
\(D^4\)\_s times \(S^3\)\_t and \(S^3\)\_s times \(D^4\)\_t, with their physical
weights (-1,-1,+1,+1). The actual contracting direction enters
the first and exits the second, exactly as in EA.40--41.
Their target-minus-source full face is therefore one positive
\(S^7\) with its charge-one physical circle and line O(-1).
Call its physical quotient \(P_0\). Its physical characteristic pair
is (-1,-4), so its index is zero. Its label is the first primitive
piece; it is not glued to the second primitive origin by a source
period. That period preserves the intermediate root character,
whereas passing from object zero to object one changes it by
m/2. The original B closes only after both pieces.

The \(q_1\) factor occurs only in \(K_1\). Use its actual associated node
and the compact domain D1 of (GC.3), with the lower horn and
full charged annulus of GC.2. The complete chain identity
EA.50--52, now with the exact GC.4--6 defining functions, gives
one further positive full charge-one link \(P_1\). Its source/line
node coefficient includes \(c_{lambda}\)/c0. There is exactly one
adjacent-node height shift from this \(q_1\) factor. It is retained
in the final B endpoint (GC.25)--(GC.27), and is not assigned
a second time to \(K_0\). The physical pair of \(P_1\) is again (-1,-4).
This chart has no isolated positive-chamber physical fixed point:
z has auxiliary weight zero, and the only fixed locus is the
whole surface z=V=0 already identified in GC.1.

The remaining upper face of this associated cut is the complete
free product of GC.5, including its exceptional \(\mathbb P^1\). Use the
proper graph and both target charts of GC.6 to carry its actual
inverse primitive to the intermediate object. Insert the finite
source trace of GC.8 on this free upper part. Its stable side
is the exact norm-strip map, because the positive auxiliary flow
preserves the common source matrix and the unit E. Its target
line coefficients are precisely \(L_1\),\(L_2\) and the composed \(F_{B}\) L.
The source Tate-node and associated-line-node arrows (GC.23b)--
(GC.23c) retain the final projective corner. Thus this upper
insertion has one additional full free face \(Q_{up}\), and does not
add another charge-one origin or another fixed surface.

The fixed locus is joined to the outside by the actual RT.1
normalization immersion and the full EA.42--46 Cartier map.
It is the whole F=\(\operatorname{Bl}(\mathbb P^1\times D)\), with its one exceptional E.
\(K_0\) is a nonzero-line piece and the inserted upper trace is free,
so neither contributes another fixed component. The local fixed
\(S^3\) is retained inside F and receives no separate \(D^4\) cap. The
one global boundary \(S^2\times S^1\) is capped by the actual C0
product RT.33a--b with its positive-spinor coefficient arrow.
Use the exact RT.34 triangle bordism to compare it with \(C=\mathbb P^1\times D^2\), retaining both g/w clutches and the whole \(S^4\) normal
bundle. This contributes the two deleted-tube faces -\(B_F\) and -\(B_4\).

Every remaining ordinary B face is paired by K1K0=B; every
root face is paired through (GC.31); and the source-base collar
is paired through the full TL \(\theta\)/Cartier map. The inequalities
\textbar{}\(q_1\)\textbar+\textbar{}\(q_2\)\textbar\textless=R place all inner changes away from the prescribed
outer parameter circle and source-base boundary. The normal
maps on intersecting collars are their written common Cox maps;
their products on a triple overlap are the actual identity.
The two chronological determinant roots and their finite
characters are those of GC.4. The free upper trace has its
specified determinant frame and physical coefficient. Thus no
unrecorded order-two spin character is used in joining these
faces. The construction uses finitely many compact norm pieces,
the two actual source projective charts at each primitive step,
and the one compact source-moment trace. All identifications are
between its specified closed collars. Rounding yields the finite
smooth oriented orbit bordism.

Let N be its actual outside conifold/Hopf completion. Give \(Q_{up}\)
the sign induced by its boundary inclusion in this construction.
The resulting physical orbit boundary has the full face count

\[
                  -N-P_0-P_1-B_F-B_4+Q_{\rm up}.
 \tag{GC.32}
\]

The minus signs of \(P_0\) and \(P_1\) follow from their positive total-space
finite-end chains and the circle-first boundary reversal. The
fixed-tube signs are the actual RT.34 cap signs. The free face
\(Q_{up}\) has its independently specified trivial physical line.
There is no charge-zero projective substitution or an additional
isolated positive-chamber point in this ledger.

The characteristic expression on \(B_F\) is the full RT.32 expression
with a=1 and both retained clutch integers, hence is zero. On \(B_4\)
it is zero with arbitrary normal \(c_2\) retained. The two charge-one
faces and \(Q_{up}\) have just been computed as zero from their actual
line maps. The outside characteristic expression is therefore
zero by Stokes for the carried physical line and stable tangent.
The outside boundary is the actual fixed-scalar source loop,
because all constructions are identity modifications off the
inner collision collar and retain (GC.31) on the cap. The
physical/Jin inverse-return arrow is the one in K68; the zero
value is unchanged by that explicitly carried reversal. This
completes the good-branch boundary proof. Section GC.10 gives
the retained-clutch evaluation and the full actual-return arrow.

\par

\medskip\noindent\textbf{GC.10. The carried characteristic expression and actual parameter return}

\par\nobreak

Write M for the actual physical nonzero-vector line and x=\(c_1\)(M).
On the free quotient bordism choose invariant connections product on
all the specified collars and patch them on its compact interior.
The circle tangent is oriented and trivial, so the actual tangent
splitting identifies its \(p_1\) with the pullback of the quotient \(p_1\).
Stokes applies to the closed polynomial

\[
                      A(Y)=\int_Y(4x^3-p_1(TY)x)/24.
 \tag{GC.33}
\]

Every line in this formula is the one carried by the written
coefficient maps. In particular the two chronological charged
faces have x=-h and \(p_1\)=4h squared on their positively oriented
\(\mathbb{CP}^3\) quotients, and hence A(\(P_0\))=A(\(P_1\))=0. The upper free face has
its actual coefficient frame L and x=0, so A(\(Q_{up}\))=0.

Here is the whole fixed-cap evaluation with the actual clutches
retained. On F-C=\(\operatorname{Bl}(\mathbb P^1\times\mathbb P^1)\), use H,B,E with H squared and
B squared zero, HB=1, E squared=-1 and HE=BE=0. If the full
normal/cap coefficient ratios have the integers w+ and w- of
RT.29, the two positive spinor line classes for a=1 are

\[
        l_+=H-E+w_+B,\qquad l_-=-H+w_-B.
\]

Thus b=\(c_1\)(\(E_{spin}\))=-E+(w\_++w\_-)B has b squared=-1 and
\(c_2\)=w\_--w\_+. The actual Hopf coefficient arrow makes M=O(-1)
on its projectivization. With h=\(c_1\)(O(1)), the relation
h squared+b h+\(c_2\)=0 and pi\_*h=1 give

\[
 k=c_2-b^2=1+w_--w_+,\qquad
 \ell=-3\sigma-b^2+4c_2=4+4(w_--w_+).
\]

Here \(\sigma\)=-1. Consequently A(\(B_F\))=0 with both integers retained.
The exact triangle cap comparison gives a normal bundle on \(S^4\)
with arbitrary retained \(c_2\); its pair is (\(c_2\),4c2), so A(\(B_4\))=0.
There is no use of a vanishing normal clutch. Substitution in
the actual oriented ledger (GC.32) therefore gives A(N)=0.

To identify this with the requested parameter return, take a
physical-\(\delta\)-averaged connection on the actual completed
family, product on the selected tube in the exact RT.33a--b
Cox/Hopf coordinates. Its conifold quotient return is a
relative diffeomorphism f of \(S^2\) times \(S^3\), with the stated
identity product on that tube. It preserves the actual Euler
class, which generates \(H^2\) cohomology. Its orientation then
preserves \(H_3\), so f is in the required homology-trivial relative
group. The physical Hopf lift of P3F.22 is literally the return
on this same complement and the same identity extension on
the selected tube. P3K.68 supplies the full oriented pieces
and the inverse-clutch arrow

\[
                  \lambda_M(f)=-\Sigma(f^{-1})=\Sigma(f).
 \tag{GC.34}
\]

Hence the fixed-scalar source return has physical integer
A(N)=0 and oriented class zero in \(\Theta_7\). This last step uses
the actual source, cap, tube and line maps; it does not select
a return representative from the zero polynomial alone.

For the literal reflected fan the already proved single
native Cox comparison of P3K.60 has free pair (1-a,a)=(0,1).
The same physical relative convention of P3K.68--69 therefore
adds one to the central class. Its reflected fixed-scalar
class is one. This does not assign a new reflected integer
characteristic pair from that residue.

Finally the canonical scalar-root loop in this good branch
has physical index zero in these two actual fans, by the
retained scalar page computation and its physical-line
correction, with the narrow native C2 comparison of P3K.62.
The exact RT.12 based phase homotopy, at
\(r_{base}\)=v\_*\^{}k \(r_0\) and \(r_0\)\^{}m=lambda0, adds k such actual scalar
turns to the fixed-scalar source loop. Thus every retained
coupled loop \(\lambda\)=lambda0 v\^{}(km) has the same two oriented
classes zero and one. Both norm signs and all realized finite
\(\delta\) characters were kept in (GC.9)--(GC.10) and (GC.31).
No class in \(\Theta_7\) has been divided.

The certificate \path{code/part3/source_zero_good_cut.py} checks the
exact cut, node, \(\theta\)-factor, cap-square, determinant and
moment identities. Its first extended run reached the literal
charge-one index and encountered a Python float-zero versus
exact SymPy zero comparison; the code and output are retained
in \path{records/attempts/source_zero_good_trace_float_zero}. The
correction uses Rational for that literal zero only. All
preceding geometric algebra had already passed. The current
71-check certificate complements the geometric face proof
and its independent reviews.
\end{proof}
